\section{Sharp isolated-solution growth at fixed derivative order}
\label{app:isolated-sharpness}

The quadratic isolated family of Appendix~\ref{app:actual-first-order}
extends to every fixed derivative order. Here we give a matching general
upper bound and an explicit family attaining its exponent. These are
bounds on the actual polynomial-solution locus, before any agreement
filter. The attaining family still has a constant nearby-label count at
a fixed positive gap, so this sharpness does not establish sharp MCA error.

\begin{theorem}[Isolated regular solution upper bound]
\label{thm:isolated-upper}
Let $d\ge1$, $D\ge d$, and let $E$ be algebraically closed of
characteristic zero or $p>D$. Let $Q(X,z,u_0,\ldots,u_d)$ have total
jet degree at most $B\ge1$ and challenge degree at most $H$, with
unrestricted $X$-degree. Let $J_0$ count the isolated points of the
reduced closure of actual regular pairs
\[
 Q(X,z,P,P',\ldots,P^{(d)})=0,\qquad \deg P\le D,\qquad
 Q_{u_d}(X,z,P,\ldots,P^{(d)})\not\equiv0.
\]
Put
\[
 \tau=\max\{0,2(D-d)-1\},\quad b=1+\tau(B-1),\quad
 K=H+B(b+\tau H),\quad q=B+H.
\]
Then $J_0\le qK^{d+1}=O_{d,B,H}(D^{d+1})$.
\end{theorem}
\begin{proof}
If $J_0=0$ there is nothing to prove. Every isolated point of this
closure belongs to the actual regular locus, since that locus is dense.
Choose a Taylor center $t$ at which the separant is nonzero at every
isolated regular solution; there are finitely many such points and
only finitely many excluded centers for each. Write
\[
 S=Q_{u_d},\qquad A_0=Q_X+\sum_{i<d}u_{i+1}Q_{u_i},\qquad
 \mathcal D=\partial_X+\sum_{i<d}u_{i+1}\partial_{u_i}
                    -(A_0/S)\partial_{u_d}.
\]
The quotient-rule recurrence in the proof of
Theorem~\ref{thm:actual-first-order-components}, with this derivation,
gives for the $j$th derivative, $j>d$, denominator
$S^{2(j-d)-1}$, jet numerator degree at most
$1+(2(j-d)-1)(B-1)$, and challenge degree at most $(2(j-d)-1)H$.
The initial derivatives $u_0,\ldots,u_d$ need no denominator.
Thus the degree-$D$ Taylor reconstruction, with common denominator
$S^\tau$, has coefficient numerators of jet degree at most $b$ and
challenge degree at most $\tau H$. Taylor factorials are invertible.

Substitute this polynomial and its first $d$ derivatives into $Q$,
clear $S^{\tau B}$, and take coefficients in $X$. Every residual has
total initial-data degree at most $K$: a term of jet degree $e\le B$
has numerator jet degree at most
$eb+\tau(B-e)(B-1)\le Bb$ and challenge degree at most
$H+\tau BH$. Include the initial equation
$Q(t,z,u_0,\ldots,u_d)=0$, of degree at most $q$.
On $S(t,z,u_0,\ldots,u_d)\ne0$, these equations describe precisely
the initial jets of actual degree-$\le D$ solutions. The initial-jet
and reconstruction maps are inverse regular maps there, so they
preserve the isolated points being counted.

The initial hypersurface has dimension at most $d+1$ in $d+2$
variables. Cut successively by $d+1$ generic linear combinations of
the residuals, each of degree at most $K$. At each stage, discard
positive-dimensional components on which every residual vanishes;
none can contain a target isolated point in the regular open locus.
On each remaining component through a target point, some residual
is nonzero, and a generic combination cuts its dimension. Choose the
combinations simultaneously for the finitely many components at
each step. This process retains every target isolated point.
B\'ezout bounds the sum of the degrees of the retained intersections
by $qK^{d+1}$. This argument does not require that the whole unpruned
intersection be zero-dimensional: other positive-dimensional
components may remain away from the target points.
\end{proof}

\begin{theorem}[Attaining every isolated-count exponent]
\label{thm:isolated-lower}
Fix $r=d+1\ge2$ and $D\ge1$, in characteristic zero or $p>D+r$.
Let $R$ be monic and squarefree of degree $N=D+1$, with nonzero roots.
Put
\[
 \phi_j=X^j\ (1\le j<r),\qquad \phi_r=X^r+z,
 \qquad F_j=R\phi_j'-P\phi_j,
\]
and let $Q_r$ be the Wronskian of $F_1,\ldots,F_r$, with derivative
rows $0,\ldots,r-1$ and $z$ held constant. Then $Q_r$ is a nonzero
differential polynomial of order $d$, jet degree $r$, challenge
degree one, and degree one in its highest derivative.
Its degree-$\le D$ polynomial solutions are exactly
\[
 P=R\frac{H'}H,\qquad z=H(0),
\]
where $H$ ranges over monic degree-$r$ polynomials whose root multisets
lie among the roots of $R$. All are regular, and there are exactly
\[
 \binom{D+r}{r}=\Theta_d(D^{d+1})
\]
isolated polynomial--challenge pairs, with no positive-dimensional
solution component. For fixed $d$, all labels can be distinct over
prime fields of size $O_d(D^{2d+1})$.
\end{theorem}
\begin{proof}
The separant has the explicit value
\begin{equation}
 \frac{\partial Q_r}{\partial P^{(r-1)}}
 =(-1)^r C_r R^{r-1}\bigl(X^r+(-1)^{r-1}z\bigr),
 \qquad C_r=\prod_{j=1}^{r-1}j!.
 \label{eq:isolated-separant}
\end{equation}
Indeed, the highest jet occurs only in the last derivative row, with
coefficient row $-(\phi_1,\ldots,\phi_r)$. Move that row to the top.
In the other rows the successive highest derivatives of the $\phi_j$
have coefficient $R$, giving $(-1)^rR^{r-1}\operatorname{Wr}(\phi_1,\ldots,\phi_r)$.
The monomial Vandermonde formula gives
\[
 \operatorname{Wr}(\phi_1,\ldots,\phi_r)
   =C_r\bigl(X^r+(-1)^{r-1}z\bigr).
\]
This is nonzero for every constant $z$. The pure $P^r$ coefficient
is $(-1)^r\operatorname{Wr}(\phi_1,\ldots,\phi_r)$, so the jet
degree is exactly $r$. All claimed degree bounds follow; the
$X$-degree is also $O_r(D)$, although the upper bound does not need it.

Each specialized $F_j$ has degree at most $D+r<p$. In this range a
polynomial Wronskian vanishes exactly when its columns are linearly
dependent over constants. To see the nontrivial direction, echelon
any independent polynomial basis to distinct leading degrees. The
leading Wronskian coefficient is their nonzero leading-coefficient
product times the Vandermonde of these distinct degrees, all less
than $p$. The same argument works in characteristic zero.

The first $r-1$ columns are independent for every polynomial $P$.
Otherwise $RH'-PH=0$ for a nonzero
$H\in\operatorname{span}(X,\ldots,X^{r-1})$. Its root at zero has
multiplicity between one and $r-1$, so $H'/H$ has a nonzero simple
pole there, contradicting $R(0)\ne0$ and polynomial $P$.
Consequently $Q_r=0$ if and only if a unique monic polynomial
\[
 H=X^r+\sum_{j=1}^{r-1}h_jX^j+z
\]
satisfies $RH'-PH=0$. Its root multiplicities are at most $r<p$;
polynomiality forces every root to belong to $R$, and conversely
every such root multiset works. Partial-fraction residues distinguish
all the multisets. The leading coefficient of $P$ is $r$, so its
degree is exactly $D$. Equation~\eqref{eq:isolated-separant} proves
regularity, and the classification proves finiteness and the count.

To separate the labels, greedily choose $N$ nonzero roots whose
products of $r$ factors, with repetitions allowed, are all different.
Suppose $k\ge1$ roots have been chosen. Products of any fixed smaller
number of factors are also distinct, by padding with one old root.
A new product has form $x^jc$, with $1\le j\le r$ and $c$ an old
product of $r-j$ factors. Its collision with an old $r$-fold product
excludes at most $j$ choices of $x$. Collisions between $x^ic$ and
$x^jc'$ for $i<j$ exclude at most $j-i$ nonzero choices. Equal-power
collisions are already excluded. Writing
$M_s(k)=\binom{k+s-1}{s}$, including $M_0(k)=1$, a bound on all
forbidden choices, including old roots and zero, is
\[
 k+1+\sum_{j=1}^r jM_{r-j}(k)M_r(k)
 +\sum_{1\le i<j\le r}(j-i)M_{r-i}(k)M_{r-j}(k)
 =O_r(k^{2r-1}).
\]
Start from one nonzero root and choose a prime larger than this bound
for every $k<N$ and larger than $D+r$. A prime of size
$O_r(N^{2r-1})$ suffices. All resulting $H$, $P$, and distinct labels
$H(0)=(-1)^r\prod a$ lie in that prime field.
\end{proof}

\begin{proposition}[The attaining family still has constant nearby count]
\label{prop:isolated-higher-agreement}
For the family in Theorem~\ref{thm:isolated-lower}, let a received line
have $L$ distinct labels with candidates agreeing on at least $A$
of $n$ coordinates, where $1\le A\le n$. Let $r_0$ count domain roots of $R$, and put
$m=n-r_0$, $t=A-r_0$, and $C=3r-2$. If $t^3>8Cm^2$, then
\[
 L\le\max\left\{\frac{4m}{t},
             \frac{8m^3(r-1)}{t^3-8Cm^2}\right\}.
\]
In particular $L=O_{r,\eta}(1)$ at fixed $A-D\ge\eta n$ and
sufficiently large $n$. This conclusion remains valid for any
assignment of one label to each candidate, counting distinct labels.
\end{proposition}
\begin{proof}
Choose one candidate at each selected label, and normalize it to
$W=P/R=H'/H$. Its residues at the roots of $R$ belong to
$\{0,\ldots,r\}$, and their sum is $r$. Distinct candidates have
different residue vectors. On an affine line of labeled functions
$(z,W)$ through two of them at distinct labels, choose a residue
coordinate where their vectors differ. That residue is an injective
affine function of $z$ along the line and has only $r+1$ allowed
values. Thus every such line contains at most $r+1$ candidates.
There are at most $(r-1)\binom L2/3$ collinear triples.

A noncollinear triple gives a nonzero rational relation
$\sum_i c_iW_i$, with $\sum_i c_i=\sum_i c_iz_i=0$.
On the product of its three degree-$r$ denominators, the numerator
has degree at most $3r-2$: each $W_i$ starts with $r/X$ at infinity,
so the highest coefficient cancels. Such a triple has at most
$C=3r-2$ common agreeing coordinates outside the roots of $R$.
A collinear triple has at most $m$.

Let $h_x$ count agreements at each outside coordinate. Each selected
candidate contributes at least $t$ agreements. The same convexity
estimate as in Proposition~\ref{prop:isolated-family-mca} gives
\[
 m\left(\frac{tL}{m}-2\right)_+^3
 \le6\sum_x\binom{h_x}{3}
 \le C L^3+m(r-1)L^2.
\]
If $L\ge4m/t$, the left side is at least $t^3L^3/(8m^2)$;
rearranging proves the second displayed alternative. Otherwise the
first alternative holds. Finally $r_0\le D+1$ gives
$t\ge\eta n-1$ and $m\le n$, proving the constant fixed-gap bound.
\end{proof}

\paragraph{Collinear triples cannot simply be omitted.}
For $r=3$, take two roots $a=1$, $b=-2$ and
$H_j=(X-a)^{3-j}(X-b)^j$ for $j=0,1,3$. Their logarithmic
derivatives are affine in $j$, and their labels $-1,2,8$ are also
affine at these three indices. Over sufficiently large characteristic
this is a distinct-label collinear triple. The extra $L^2$ term in
the proof handles such configurations.

\paragraph{Verification and limits.}
The package \path{research/isolated_solution_sharpness/} checks the
separant through order five, exhausts $760$ polynomials covering
$5{,}270$ polynomial--challenge pairs at orders one through three,
and includes characteristic-three and characteristic-five negative
controls. Prime-field examples have $969$ distinct isolated labels
at order two, $495$ at order three, and $462$ at order four.
A separate checker verifies the reconstruction recurrence on four
equations of orders two through five and $360$ finite Taylor fixtures.
The checks support the written proofs; they do not replace the generic
intersection argument. No fixed-gap superlinear MCA lower bound follows:
the polynomially many isolated solutions in the sharpness family fail
the required nearby-label condition.

\subsection{A rational envelope has only linearly many bad labels}
\label{sec:rational-envelope}

The agreement obstruction above extends beyond a fixed residue alphabet.
The following elementary incidence bound allows arbitrarily many candidates
and arbitrary characteristics. It gives another reason why a large
algebraic solution count need not make the proximity bound sharp.

\begin{proposition}[Rational-envelope bound]
\label{prop:rational-envelope}
Let $\mathcal D$ consist of $n$ distinct points of a field $F$, and fix
polynomials $S,R\in F[X]$ with $R\ne0$. Consider any family of polynomials
$P$ of degree at most $D$ satisfying
\[
 \frac{P-S}{R}=\frac{a_P}{b_P},\qquad
 \deg a_P,\deg b_P\le r,\quad b_P\ne0.
\]
Let $h$ be the number of zeros of $R$ on $\mathcal D$. For an integer
agreement threshold $1\le A\le n$, suppose
\[
 A-h\ge\beta n>0,\qquad n\ge48r/\beta^3.
\]
On every received line $f+zg$, the number of full-support MCA-bad labels
witnessed by this family is at most
\[
 \frac{16(n-A+1)}{\beta^3}.
\]
Here badness means that $g$ has no degree-at-most-$D$ interpolant on the
candidate's entire agreement set. There is no characteristic hypothesis.
\end{proposition}
\begin{proof}
First, any polynomial pencil $U+zV$, with $\deg U,\deg V\le D$, has
at most $M=n-A+1$ bad labels. Let $c$ count its persistent agreements,
where $(f,g)=(U,V)$. Each other coordinate agrees at at most one label.
If $c<A$, the number of nearby labels is at most
$\lfloor(n-c)/(A-c)\rfloor\le n-A+1$.
If $c\ge A$, a bad label must have a nonpersistent agreement, giving
at most $n-c\le n-A$ labels.

Choose one bad candidate $P_z$ at each of $L$ distinct labels and put
$W_z=(P_z-S)/R$. Delete the $h$ zeros of $R$ and normalize the received
line by subtracting $S$ and dividing by $R$. Each candidate retains at
least $\beta n$ agreements. Its reduced denominator is nonzero at every
retained coordinate, since $W_z=(P_z-S)/R$ is regular there.

A noncollinear triple of graph points $(z,W_z)$ over $F(X)$ has a
nonzero affine-incidence determinant. Clearing its three denominators
gives a polynomial of degree at most $3r$, so the triple has at most
$3r$ common retained agreements. A graph line through two selected
points corresponds, on multiplying by $R$ and adding $S$, to a
polynomial pencil of degree at most $D$: interpolate the two actual
polynomials at their distinct scalar labels. It therefore contains at
most $M$ selected bad labels.

Let $b_x$ count retained agreements at coordinate $x$, and set $b_x=0$
at deleted coordinates. If $L\ge4/\beta$, convexity gives
\[
 \sum_x b_x(b_x-1)(b_x-2)
 \ge n(\beta L-2)^3\ge\beta^3nL^3/8.
\]
Noncollinear ordered triples contribute at most $3rL^3$. Each ordered
pair determines one graph line with at most $M$ choices of a third
selected label; these triples contribute at most $nML^2$. Consequently
\[
 \beta^3nL^3/8\le3rL^3+nML^2.
\]
Since $3r\le\beta^3n/16$, this implies $L\le16M/\beta^3$.
The case $L<4/\beta$ satisfies the same bound. Over an infinite field
apply the argument to any finite selection of bad labels.
\end{proof}

For example, if $\deg R\le D+c$ and $A-D\ge\eta n$, one can take
$\beta=\eta-c/n$ when positive. For fixed $c,r,\eta>0$ the bound is
therefore $O_\eta(n)$. It applies to $P=RH'/H$ with $\deg H\le r$
without a restriction on the residues, and even permits
$r\le\beta^3n/48$. A bounded union of such families still has a linear
bound. The hypothesis does not hold for arbitrary positive-rate
candidates merely by setting $R=1$, since then the necessary residual
degree $r$ is generally too large.

The checker \path{research/prime_field_tightness/check_rational_envelope.py}
verifies $220$ rational triple determinants over $\F_{101}$ and all
full agreement supports in a pencil attaining $n-A+1=13$ bad labels.
These checks supplement the incidence proof.

\subsection{Many nearby labels force a rational path to be affine}
\label{sec:rational-path-rigidity}

There is a stronger restriction when the candidate itself is a rational
function of the challenge of bounded degree. No bound on the coefficient
heights in $X$ is needed.

\begin{theorem}[Rational-path rigidity]
\label{thm:rational-path-rigidity}
Let $D<A\le n$ be integers, and let $f+Zg$ be a received line on $n$
distinct points of a field $F$. Suppose $T(X,Z)\in F(X)(Z)$ has numerator
and denominator degrees in $Z$ at most $h\ge1$. Either
\[
 T(X,Z)=U(X)+ZV(X),\qquad \deg U,\deg V\le D,
\]
or the number of scalar labels $z$ at which $T(X,z)$ is defined, is a
polynomial of degree at most $D$, and has at least $A$ agreements with
$f+zg$ is at most
\[
 \frac{(8h+10)n}{A-D}.
\]
In the affine case there are at most $n-A+1$ full-support MCA-bad labels.
These statements hold in every characteristic.
\end{theorem}
\begin{proof}
Write $T=a/b$ with coprime $a,b\in F[X,Z]$, $b\ne0$, and both
$Z$-degrees at most $h$. For every polynomial specialization $P_z$ under
consideration, $a(X,z)=b(X,z)P_z(X)$ as polynomials. Put
\[
 F_x(Z)=a(x,Z)-b(x,Z)(f(x)+Zg(x)).
\]
Call a coordinate persistent if $F_x\equiv0$, and let $H$ count them.
At a persistent coordinate $b(x,Z)\not\equiv0$, since otherwise $X-x$
would divide both $a$ and $b$. Thus all but at most $h$ selected labels
agree there. At any other coordinate agreement requires a root of the
nonzero polynomial $F_x$, of degree at most $h+1$. This argument includes
points where specialization makes both $a(x,z)$ and $b(x,z)$ zero.

Put $\delta=A-D$. If $H\le D+\delta/2$, each nearby label needs at
least $\delta/2$ nonpersistent agreements. Their total incidence is at
most $(h+1)n$, giving at most $2(h+1)n/\delta$ nearby labels.

Now suppose $H>D+\delta/2$, and choose $L$ nearby labels. If $T$ is
affine in $Z$ over $F(X)$, two polynomial values give the stated
polynomial affine form; the case $L\le1$ already satisfies the numerical
bound. Otherwise any graph line over $F(X)$ meets the graph of $T$ in
at most $h+1$ labels: the numerator of $T-U-ZV$ is nonzero and has
$Z$-degree at most $h+1$.

The selected graph points $(z,P_z)$ therefore have at most
$(h-1)L(L-1)$ ordered collinear triples. A noncollinear triple has a
nonzero affine-incidence determinant, a polynomial in $X$ of degree
at most $D$, so it shares at most $D$ persistent agreeing coordinates.
Every persistent coordinate has at least $L-h$ agreements. For
$L\ge h+2$, counting ordered triples gives
\[
 H(L-h-2)^3\le DL^3+H(h-1)L^2.
\]
Using $(L-h-2)^3\ge L^3-3(h+2)L^2$ yields
\[
 (H-D)L\le(4h+5)H,
 \qquad L\le\frac{(4h+5)H}{H-D}
       <\frac{(8h+10)n}{\delta}.
\]
The smaller-$L$ case satisfies the same bound. In the affine alternative,
the sharp polynomial-pencil bound from the proof of
Proposition~\ref{prop:rational-envelope} gives $n-A+1$ bad labels.
\end{proof}

\begin{corollary}[Degree required by a rational witness formula]
\label{cor:rational-witness-degree}
Suppose a rational function $T(X,Z)$ selects a degree-at-most-$D$ nearby
witness at each of $L>n-A+1$ distinct full-support MCA-bad labels, with
agreement threshold $A>D$. If its numerator and denominator degrees in
$Z$ are at most $h$, then
\[
 h\ge\frac{(A-D)L}{8n}-\frac54.
\]
\end{corollary}
\begin{proof}
The affine alternative can cover at most $n-A+1$ of the selected bad
labels. Rearrange the other alternative of
Theorem~\ref{thm:rational-path-rigidity}. A degree-zero formula is a
constant polynomial pencil whenever it has a polynomial value, so it
cannot cover the hypothesized set either.
\end{proof}

For the complete-coverage examples, $L=p-1$ and
$A-D=\eta n+1$. Any rational formula selecting witnesses throughout
the nearby set consequently has challenge degree $\Omega(\eta p)$;
when $\eta=\Theta(1/\log p)$ this is $\Omega(p/\log p)$.
This quantifies a failure of low-degree witness coherence. It is a
degree bound, not a computational lower bound: efficient witness recovery
may use branching rather than one rational formula.

The reciprocal-gap dependence in the rigidity theorem is necessary up
to constants for fixed $h$. The construction and proof in
\path{research/prime_field_tightness/RATIONAL_PATH_RIGIDITY.md} give a
non-affine projective path with order $1/\eta$ nearby bad labels at fixed
rate. The exact checker considers lengths $160,320,640$, all at rate
$1/4$ and capacity gap $1/8$: the path has respectively $41,81,161$
polynomial members, but exactly ten nearby labels, all full-support bad.
This count is for the specified path, not for all codewords on the line.

